% SOURCE_STATEMENT: "Do sequences of manifolds with nonnegative scalar curvature and no closed interior minimal surfaces other than their boundaries avoid disappearing in the limit unless their volumes tend to zero?"
% Preparation and final self-audit: CONTRIBUTING.md prompts C and D; source/audit artifacts retained in the local physics_statements_2026_10_11 batch.
% !TeX program = lualatex
\documentclass[11pt]{article}
\usepackage[margin=25mm]{geometry}
\usepackage{fontspec}
\setmainfont{Latin Modern Roman}
\usepackage{amsmath,amssymb,mathtools}
\usepackage{unicode-math}
\setmathfont{Latin Modern Math}
\usepackage{xurl,hyperref}
\hypersetup{colorlinks=true,urlcolor=blue}
\setcounter{secnumdepth}{0}
\setlength{\parindent}{0pt}
\setlength{\parskip}{5pt}
\setlength{\emergencystretch}{3em}
\newcommand{\sref}[2]{\hyperlink{source:#1:#2}{[S#2]}}
\newcommand{\cataloguescope}[1]{\paragraph{Recorded target.}#1}
\newcommand{\R}{\mathbb R}
\newcommand{\C}{\mathbb C}
\newcommand{\N}{\mathbb N}
\newcommand{\Z}{\mathbb Z}
\newcommand{\Q}{\mathbb Q}
\newcommand{\F}{\mathbb F}
\newcommand{\eps}{\varepsilon}
\begin{document}
\section{Non-Disappearance of Nonnegative Scalar-Curvature Limits}\label{statement}

\subsection{Definitions and mathematical statement}

\paragraph{Shared notation.}$\mathbb N=\{1,2,\ldots\}$; the standard number fields have their usual meanings. All additional conventions are stated locally.


\paragraph{Definitions and conventions.}
Scalar curvature is the trace $R_g$ of the Ricci tensor of a Riemannian metric. A minimal hypersurface has mean curvature zero; a closed interior minimal surface is compact, boundaryless and lies in the interior. In a geometric convergence theory, disappearance means that a region with nonzero volume has no surviving part in the limiting geometric object. Volume is the Riemannian measure of the manifolds in the sequence. The report uses intrinsic flat convergence of oriented integral current spaces; disappearance of the whole sequence means convergence to the zero current space, $d_{\mathcal F}(M_j,0)\to0$. The cited conjecture asks whether this can happen with volumes bounded away from zero under the stated curvature and minimal-surface conditions.

\paragraph{Target retained from the canonical source catalogue.}
\paragraph{Catalogue wording retained for traceability.}
Do sequences of manifolds with nonnegative scalar curvature and no closed interior minimal surfaces other than their boundaries avoid disappearing in the limit unless their volumes tend to zero?

\paragraph{Source scope and qualification.}
The convergence topology was recovered as intrinsic flat convergence (PDF page 50). The source reports non-disappearance under additional nonnegative-Ricci-curvature or uniform-contractibility assumptions; these are partial cases, not replacements for the scalar-curvature target.
\subsection{Short English statement}
Can positive-volume regions disappear in limits of nonnegative-scalar-curvature manifolds without interior minimal surfaces?
\subsection{Sources}
\begin{itemize}
\item[\hypertarget{source:statement:1}{S1}] ULAM.AI and the UnsolvedMath Contributors, UnsolvedMath, catalogue entry. Catalogue curation under CC BY 4.0; definitions and English exposition added locally.
\url{https://huggingface.co/datasets/ulamai/UnsolvedMath}.
\item[\hypertarget{source:statement:2}{S2}] Mathematical Aspects of General Relativity (2012), Oberwolfach Reports, DOI 10.4171/owr/2012/37, pp. 2269--2333. Primary report accessed 2026-10-11; relevant PDF pages 50, 51.
\url{https://doi.org/10.4171/owr/2012/37}.
\item[\hypertarget{source:statement:3}{S3}] Sormani, scalar curvature and intrinsic-flat convergence (catalogue-listed context; not independently consulted in this preparation).
\url{https://arxiv.org/abs/1606.08949}.
\end{itemize}
\label{end:statement}
\end{document}
